\documentclass[12pt]{article}
\usepackage[utf8]{inputenc}
\usepackage[LGR, T1]{fontenc}
\usepackage[greek.ancient,english]{babel}
\usepackage{alphabeta}

\usepackage[letterpaper, margin = 1in]{geometry}
\usepackage{fancyhdr}
\usepackage{titling}
\usepackage{titlesec}
\usepackage{hyperref}
\usepackage{lastpage}
\usepackage{amssymb}
\usepackage{setspace}
\usepackage{tabularx}
\usepackage{array}
\renewcommand{\arraystretch}{1.5}

\makeatletter
\newcommand{\hangpara}[2]{\hangindent#1\hangafter#2\noindent}
\newenvironment{hangparas}[2]{\setlength{\parindent}{\z@}%
  \everypar={\hangpara{#1}{#2}}}{\par}
\makeatother

\providecommand{\medium}{\normalsize}

\title{Commodity Value in Marx and Aristotle}
\author{Ian Thomas White}
\date{}

\makeatletter
\renewcommand{\maketitle}
{
\begin{flushleft}
\textbf{{\large\@title}}
\newline
\@author
\end{flushleft}
}
\makeatother

\titleformat*{\section}{\large\bfseries}
\titleformat*{\subsection}{\medium\bfseries}
\titleformat*{\subsubsection}{\small\bfseries}

\pagestyle{fancy}
\fancyhf{}
\setlength{\headheight}{14.5pt}
\lhead{Ian Thomas White}
\chead{Commodity Value in Marx and Aristotle}
\rhead{\thepage \ of \pageref{LastPage}}

\doublespacing
\widowpenalty=100000
\clubpenalty=100000

\begin{document}

\thispagestyle{empty}

\begin{singlespacing}
\maketitle
\tableofcontents
\end{singlespacing}

\newpage

\section{Introduction}
I'm going to attempt to explicate something hopefully not too ``slight in content'',\footnote{Karl Marx, \textit{Capital Volume I}, trans. Ben Fowkes (New York: Penguin Books, 1990), 90.} namely, a small passage in \textit{Capital} in which Marx references Aristotle. With recourse to Marx's implicit quantity-quality characterizations of homogenized labor, I hope to show how a \textit{prima facie} account of Marx's incorrect reading of the Aristotle passage might be flipped aright.\footnote{In this paper, unless otherwise marked, references to Aristotle are my own translations of the Greek texts (in the cited Loeb editions). References to Marx are generally to the English translations directly, but for the German source I used the online version of \textit{Das Kapital}: Karl Marx, \textit{Das Kapital Band I}, Karl Marx-Friedrich Engels-Werke, Band 23 (Berlin: Dietz Verlag, 1968),  \url{http://www.mlwerke.de/me/me23/me23\_000.htm}.}

\section{Problema}
I must first act the pedant in order to later attempt substantial fruits. Early on in \textit{Capital}, Marx invokes Aristotle, via the \textit{Nicomachean Ethics}, as a particularly prescient forerunner to his own commodity form and as, thereby, a forerunner to the sticky value inherent therein. Aristotle, Marx says, was on the cusp of discovering the concept of value, but the former failed---but for the fact of that slave society in which he found himself. Marx first claims that Aristotle understood that the money-form stems from the simple form of value equality. Marx proceeds:  

\begin{quote}
\begin{singlespace}
``There can be no exchange,'' he [Aristotle] says, ``without equality, and no equality without commensurability [\textit{Kommensurabilität}]'' (``οὔτ᾿ ἰσότης μὴ οὔσης συμμετρίας''). Here, however, he falters, and abandons the further analysis of the form of value. ``It is, however, in reality, impossible (``τῇ μὲν οὖν ἀληθείᾳ'') that such unlike things can be commensurable'', i.e. qualitatively equal. This form of equation can only be something foreign to the true nature of the things, it is therefore only ``a makeshift for practical purposes'' \dots  What is the homogeneous element, i.e. the common substance [\textit{Substanz}]\footnote{Importantly, the German words have similar connotations as their English equivalents.} [?] \dots  Such a thing, in truth, cannot exist, says Aristotle.\footnote{Marx, \textit{Capital}, 151; Marx, \textit{Kapital}, 74; Greek in original, cf. Aristotle, \textit{The Nicomachean Ethics}, trans. H. Rackham, Loeb Classical Library 73 (Cambridge, MA: Harvard University Press, 1934), 286.}
\end{singlespace}
\end{quote}

Marx continues to claim that the sought-after ``homogenous element'' is human labor. This ``commensurability'' couldn't, for Aristotle, occur ``in reality'' because his slave society barred the concept of (widespread) human equality, which in turn barred the conception of a substance composed of such equalized flesh. In line perhaps with Marx's here likely-intentional blustery presentation, his argument seems to fall flat for a number of reasons. I confess, though, that I will, eventually, attempt to restitute the passage.   

Marx's paraphrase of \textit{Kommensurabilität} for συμμετρίας as ``qualitatively equal'' is unwarranted, and, if anything, the term emphasizes quantitative aspects over qualitative ones. Marx's stress on the commonplace dative adverbial ``in reality'', here hedged by particles, is similarly too strong. The passage quoted by Marx, furthermore, is part of an ethical argument in which Aristotle wants to define (one sense of) justice as a mean between too much and too little.\footnote{Whether this mean-based account works out for Aristotle's sense of justice is another matter.} So where Marx reads Aristotle's ``money'' (νόμισμα\footnote{Ibid., 182.} [lit. customary usage])- equals-bed claim as something hinting at (some kind of separately existing) substance, Aristotle is only trying to describe a measuring tool for a particular subset of a particular type of justice.\footnote{Namely, the ``corrective'' (διορθωτικόν) form. Ibid., 172.} The imputation by Marx of Aristotle's substance search is thereby equally unwarranted. It's particularly telling that, in almost the same breath, Aristotle describes both an actual human judge and money by the same term, namely as a ``mean'' or ``middle (term)'' (μέσος), and he claims that ``if people gain the mean [here, a judge], they will gain justice''.\footnote{Ibid., 276, 282.} Finally, Aristotle's emphasis on proportionality (usually a form of ἀναλογ-)\footnote{E.g., ibid., 280.} between two parties as the arbiter in matters of justice, along with other related \textit{sym}\- (συμ-) words, makes ``proportionality'' a better translation for συμμετρίας than ``commensurability'', at least insofar as ``commensurability'' might carry the unwanted fat of some proper substance. Here's my translation:  

\begin{quote}
\begin{singlespace}
Thus money [lit. customary usage], acting just like a proportionate measure, makes [these objects] equal. For there could neither be association without there being exchange, nor exchange without equality, nor equality without proportionality. While, then, it may in reality be impossible for things differing so much to become proportionate, it is sufficiently possible according to our demand [lit. need].\footnote{Ibid., 286.}
\end{singlespace}
\end{quote}

This ``sufficiently possible \dots  demand'', what Marx calls ``a makeshift for practical purposes'', at the end, is, moreover, what Aristotle previously describes as the basis for money qua custom, and it acts not as a mere repair strategy but as an explicit pragmatic measure.\footnote{Ibid., 284.} More damning for Marx's ultimate \textit{ex causa servorum} is that Aristotle introduces the at least putatory, in exchange, equality between unequal persons:   

\begin{quote}
\begin{singlespace}
For association doesn't occur between two doctors, but between a doctor and a farmer, even between those completely different and not equal. But these same must become equal to each other.\footnote{Ibid., 282.}
\end{singlespace}
\end{quote}

\noindent And somewhat later:

\begin{quote}
\begin{singlespace}
There, then, will be reciprocity, whenever [certain goods] are equalized, so that, just as the farmer (is) to the shoemaker, so the work of the shoemaker (will be) to that of the farmer.\footnote{Ibid., 284.}
\end{singlespace}
\end{quote}

\noindent It seems difficult to claim that a lack of widespread equality caused Aristotle's deficiency when Aristotle took person-equality as both a commonplace occurrence and as an essential procedural step in his own ethical apparatus. Marx himself, in another poorly translated passage, notes Aristotle's ``chrematistics'', as arising specifically to deal with (widespread) money interactions, so Marx's equality-permanence claim becomes another unjustified grab at an unfound substance in Aristotle. Nor, I'll add, do I think my pedantry would be well denied by mentioning Aristotle's dehumanized slave from the \textit{Politics}\footnote{Aristotle, \textit{Politics}, trans. H. Rackham, Loeb Classical Library 264 (Cambridge, MA: Harvard University Press, 1934), 14-16.} as something that displays a qualitative difference. Since that slave is strictly property, no different from Marx's disintegrating needle unto the yarn's warp and woof, the slave thereby wouldn't count as labor, anyway---i.e., the unequal farmer and doctor, as laborers, form the operative set of labor relations. Still, Aristotle did fail to stumble upon coagulated labor-time, and it may be fitful to find a reason.\footnote{I don't mean to deny that certain sociological conditions constrain potential conceptual insight, only that Marx's such account is here inadequate, which then suggests the need for other considerations.}

\section{Residua}
Before I root out Aristotle's deficiency, I'll take a look at Marx's depiction of commodity value in \textit{Capital}. Marx primarily uses two contrasting images to describe the mass of labor-time that constitutes this value: that of a gelatinous substance called \textit{Gallerte}, usually translated with some form of ``congealed'', and that of a crystal (\textit{Kristall}). I say ``contrasting'' because, though they both varyingly refer to the same amassed abstract human labor, their emphases as images are inverses. \textit{Gallerte}, rather than an abstract noun, is a ``specific food product, an aspic of jelly made from the ground-up parts of animals''\footnote{Sianne Ngai, \textit{Our Aesthetic Categories: Zany, Cute, Interesting} (Cambridge, MA: Harvard UP, 2012), 67. Ngai gets this characterization from Keston Sutherland.}---a gelatin. As gelatin, it would appear as a kind of surfaceless blob, with a body such that where ``it'' begins is difficult to parse. Always enterable, were one to softly poke it, one would obtain the squish of lost finger. Alternately, a crystal, an appropriately abstracter term, would appear as if solely surface, sleekly closed from prying nails.  

I don't think it would be farfetched, then, to read the pair of terms as the depiction of a quantitative and qualitative distinction that inheres within abstracted labor itself---two ways to view abstracted labor, in a book that is, of course, well-invested in this sort of binary. The distribution of these descriptors suggest support for the reading. \textit{Gallerte} shows up only in \textit{Capital} chapter one, ``The Commodity'' chapter, when this concept is still new and difficult, and when otherwise previously-observed phenomena must be reconstituted accordingly.\footnote{``\textit{With the exception of the section on the form of value}, therefore, this volume cannot stand accused on the score of difficulty.'', Marx, \textit{Capital}, 90, emphasis added. A chapter in which, I'll add, even Marx at one point in desperation abruptly claims ``the problem is already solved'' and swiftly continues on, ibid., 142.} \textit{Kristall}, though, shows up mostly in chapters two and three, on exchange and circulation, as well as sporadically after that throughout the book.\footnote{This distribution may be what tacitly convinced Fowkes to translate another term, \textit{festgeronnener}, as ``congealed'' in chapter one and as ``crystallized'' in chapter seven, Ibid. 130, 297; Marx, \textit{Kapital}, 54, 204. See the attached Appendix 2 for a list of the uses of the terms.}   

More suggestively, \textit{Gallerte} shows up exclusively as a modifier of ``labor'' (\textit{Arbeit}), including a compound of the two, ``jelly-labor'' (\textit{Arbeitsgallerten}),\footnote{Marx, \textit{Capital}, 135; Marx, \textit{Kapital}, 59.} while \textit{Kristall} mostly modifies value (\textit{Wert}) itself or money/gold (\textit{Geld}), including, again, several instances of the compound, ``money-crystal'' (\textit{Geldkristall}).\footnote{Marx, \textit{Capital}, 181; Marx, \textit{Kapital}, 101.} In other words, the more gelatinous view is associated with the quantitative matter, the substance, while the crystallized view is associated with the form by which this commodity value confronts itself in exchange (where surfaces would allow a more antagonistic front). This within-value quant-qual split, to further the argument, can help formalize how ``despite its [the coat's] buttoned-up appearance, the linen recognizes in it a splendid kindred soul''\footnote{Marx, \textit{Capital}, 144.}; or, namely, this split displays the mechanism by which ``[in] the value-relation \dots  the natural form of commodity B becomes the value-form of commodity A, in other words the body of commodity B becomes a value mirror of commodity A''.\footnote{My own more literal translation, cf., ibid., 145. Marx, \textit{Kapital}, 67.} However the ambiguous possessive relationship of the mirror gets worked out, the crystallized surface of homogenized labor seems to retain its reflective properties during exchange---a within-value analogue to exchange value qua quantity and use-value qua quality, though the analogue is directly and fittingly reversed. Commodities, it seems, communicate via a chiastic structure comprised of their inner and outer quantitative and qualitative components.  

The first time both terms show up, though, they show up together, in a dripping passage, in which the contrasting view of \textit{Gallerte} and \textit{Kristall}, in their unmixed noun forms, is reproduced in a pair of \textit{hapax legomena} of \textit{Capital}. Thus, we also get the ``residue'' (\textit{Residuum}) of labor products that retain only a ``ghost-like objectivity'' (\textit{gespenstige)}---on the one hand, a sticky finger trap and, on the other, a phantom of so much surface one could walk through it.   

\begin{quote}
\begin{singlespace}
Let us now look at the residue (\textit{Residuum}) of the products of labour. There is nothing left of them in each case but the same ghost-like (\textit{gespenstige}) objectivity; they are merely gelatin (\textit{Gallerte}) of homogeneous human labour, i.e. of human labour-power expended without regard to the form of its expenditure. \dots  As crystals (\textit{Kristalle}) of this social substance, which is common to them all, they are values---commodity values.\footnote{A slightly modified translation. Marx, \textit{Capital}, 128. Marx, \textit{Kapital}, 52.}
\end{singlespace}
\end{quote}

All of this is not to say that there is some progression from substance to surface, quantity to quality. After all, when Marx says ``the bones of the saints can't withstand it''\footnote{Marx, \textit{Capital}, 229.} he invokes crystal and not bone-made gelatin. The \textit{Grundrisse}, too, indirectly warns against separating the two views too sharply, via Marx's lambasting of his forefathers for separating the coffee from the pot.\footnote{Karl Marx, \textit{Grundrisse}, trans. Martin Nicolaus (New York: Penguin Books, 1993), 687.} Again, I'm only forefronting a formal, abstracted, synchronic quant-qual distinction that holds within the value form itself, and which seems to correspond to the quant-qual distinction between use-and-exchange value.\footnote{Cf., Marx, \textit{Capital}, 128.} It's important, though, at least for my trajectory here, to remember that these congelations are still ``homogenous''\footnote{Ibid., 135.} and are, notably, entities which ``can no longer be distinguished'' (lit. ``they no longer differ (from) each other (\textit{sie unterschiedenen sich nicht})'').\footnote{Ibid., 128. Marx, \textit{Kapital}, 52.} Marx is presenting, then, a conceptual entity, homogenized labor, that is both quantized, as he explicitly claims, but that is also at least formally delimitable---insofar as a surface is imputed to it---despite being pulped to the point of indistinguishability.

\section{Primes}
I fear the foregoing was somehow both condemnably obvious and laughably obtuse, but I hope my point will stand in better relief by bringing back Aristotle, this time his \textit{Metaphysics}. Marx's abstract katabasis unto labor's residue mirrors fairly well an argument progression in Aristotle's \textit{Metaphysics}' big chapter \textit{Zeta} (VII). Therein, Aristotle, in trying to fix what he means by ``substance'' (οὐσία) picks the fourth of his given four options for such a substance---a conceptually distinguishable entity often translated as ``substrate'' (ὑποκείμενον). ``The substrate'', moreover, is ``that according to which (all) the other things are reckoned, while it itself isn't further reckoned according to anything else''.\footnote{Aristotle, \textit{The Metaphysics I-IX}, trans. Hugh Tredennick, Loeb Classical Library 271 (Cambridge, MA: Harvard University Press, 1933), 316.} Aristotle begins a proof by assuming, though it's unclear, that substance is a substrate of the substrate---``[substance] (is) that which is not reckoned according to substrate but according to that which everything (else) is reckoned.''\footnote{Ibid.} Aristotle then starts stripping away various qualities in order to arrive at his hypothesized substratum and until only some pure matter remains. Whether this stripping, though, is ``from'' something is also unclear.  

\begin{quote}
\begin{singlespace}
For when (all) other things have been stripped, nothing (else) remains. For while (all) other things are accidents, products, and powers of their bodies; length, breadth, and depth are certain quantities but not substances. For a quantity is not substance, but rather that to which these things primarily belong is substance. But, then, length and width and depth being stripped, we see nothing remaining, save for if something is (still) delimited by these things; so that matter, when being so considered, must appear (to be) the only substance.\footnote{Ibid.}
\end{singlespace}
\end{quote}

After methodically getting rid of qualities and quantities, Aristotle concludes:   

\begin{quote}
\begin{singlespace}
Thus, the final/ultimate [substrate] is of itself neither ``something'' nor a quantity nor anything else. \dots  It follows, therefore, that matter is substance. But this is impossible. For separateness and that ``something'' seem to belong to substance.\footnote{Ibid., 318.}
\end{singlespace}
\end{quote}

I don't want to get into muck about infamous ``prime matter'', but I do want to highlight the progression of the argument. Aristotle starts with an entity assumed to be a predicate of all things equally, and, in the process of deriving such an entity, he is led to the conclusion that nothing, at least in a certain sense, is left. At the very least, this something-nothing won't work for Aristotle insofar as he wants to define substance. In any case, what's most important for my purposes is that the assumption of this conceptual, hidden, homogenous substratum leads, for Aristotle, to it being defined as something lacking both quantity and quality.  

Marx, too starts a proof-like progression with the assumption of a conceptually distinguishable entity with a universal predicability. Marx claims:  

\begin{quote}
\begin{singlespace}
Therefore x boot-polish, y silk, z gold, etc., must, as exchange-values, be mutually replaceable or of identical magnitude. It follows from this that, firstly, the valid exchange-values of a particular commodity express something equal, and secondly, exchange-value cannot be anything other than the mode of expression, the 'form of appearance', of a content distinguishable from it. \footnote{Marx, \textit{Capital}, 127.}
\end{singlespace}
\end{quote}

The geometrical certainty by which Marx claims that this ``content'' is in fact distinguishable from the ``mode of expression'', is hardly demonstrated by the claim's subsequent recapitulation in equation format or by the rather red herring geometrically-based example.\footnote{Ibid. The stakes of Marx's claim here are that the (political, social) commodity domain functions in a certain way, namely that equality $\rightarrow$ substance of equality. A (potentially) valid example from a different domain doesn't, then, prove the stakes. I'm not saying Marx's conclusion is wrong, only that it is not (here) demonstrated.} Like Aristotle, then, the entity is assumed. I take Marx's ``etc.'' here gravely, as the statement of the universal predicability of this content---universal, at least, within the domain in which the exchange process occurs, the domain of commodity phenomena. Delimiting this domain, too, seems sufficient for eliding the potential difference between Aristotle's substratum hypothesis, which might otherwise be characterized as more or actually universal (and therefore as non-trivially incomparable), insofar as Marx's entity is, in fact, unlike Aristotle's, itself predicated by something, namely labor. But the commodity world is a closed set, and only commodities can ``understand the language of commodities''.\footnote{Ibid., 143.} Continuing like Aristotle, Marx, in order to arrive at this putative ``third thing'',\footnote{Ibid.} starts stripping, through abstraction, the starting non-universal entities of their particular qualities:  

\begin{quote}
\begin{singlespace}
\dots  we abstract also from the material constituents and forms which make it a use-value. \dots  All its sensuous characteristics are extinguished. \dots  With the disappearance of the useful character of the products of labour, the useful character of the kinds of labour embodied in them also disappears; this in turn entails the disappearance of the different concrete forms of labour. They can no longer be distinguished, but are all together reduced to the same kind of labour, human labour in the abstract.\footnote{Ibid., 128.}
\end{singlespace}
\end{quote}


\noindent Immediately after this, then, Marx introduces the residue and ghost-like passage from above, the first characterization of value as the homogenous, indistinguishable thing, this substratum.\footnote{Which Marx also calls it (\textit{Substrat}): Ibid., 133, 293. Marx, \textit{Kapital}, 57, 200.} Unlike Aristotle, Marx's initial universal predicate assumption has led him to an affirmation of that assumption and to a depiction of this substratum as something that maintains quantity (i.e., \textit{Gallerte}) and quality (i.e., \textit{Kristall}). It's not clear to me to what degree Marx is directly drawing on Aristotle's \textit{Metaphysics},\footnote{Though he likely would have known of the passage.} but the fact that Marx replaces the initial definition with its final derivation would, to put it suggestively, mesh well with the difference between Aristotle's more genera-based metaphysics and Marx's dialectical method---supporting claims to methodological innovation on Marx's part.\footnote{I'm thinking of Althusser here.} The analogous progressions can be schematized as below.


\begin{singlespacing}
\begin{table}[htbp]
\centering
\caption{\textbf{Substratum Argument Progressions}}
\footnotesize
\begin{tabularx}{\textwidth}{l|X|X}
%\hline
\multicolumn{3}{c}{} \\
& \textbf{Marx, \textit{Capital}, 127--128} & \textbf{Aristotle, \textit{Metaphysics}, 317--318} \\
\hline
\textbf{Definition} & [below] & The substrate [\foreignlanguage{greek}{ὑποκείμενον}] is that according to which (all) the other things are reckoned, while it itself isn't further reckoned according to anything else. \\
\hline
\textbf{Initial Assumption} & Therefore $x$ boot-polish, $y$ silk, $z$ gold, etc., must, as exchange-values, be mutually replaceable or of identical magnitude. It follows from this that, firstly, the valid exchange-values of a particular commodity express something equal, and secondly, exchange-value cannot be anything other than the mode of expression, the `form of appearance', of a content distinguishable from it. & [Substance [\foreignlanguage{greek}{οὐσία}]] (is) that which is not reckoned according to substrate but according to that which everything (else) is reckoned. \\
\hline
\textbf{Procedure} & \dots\ we abstract also from the material constituents and forms which make it a use-value. \dots\ All its sensuous characteristics are extinguished. \dots\ With the disappearance of the useful character of the products of labour, the useful character of the kinds of labour embodied in them also disappears; this in turn entails the disappearance of the different concrete forms of labour. & For when (all) other things have been stripped, nothing (else) remains. For while (all) other things are accidents, products, and powers of their bodies; length, breadth, and depth are certain quantities but not substances. For a quantity is not substance, but rather that to which these things primarily belong is substance. But, then, length and width and depth being stripped, we see nothing remaining, save for if something is (still) delimited by these things; so that matter, when being so considered, must appear (to be) the only substance. \\
\hline
\textbf{Conclusion} & They can no longer be distinguished, but are all together reduced to the same kind of labour, human labour in the abstract. & Thus, the final/ultimate [substrate] is of itself neither ``something'' nor a quantity nor anything else. \dots\ It follows, therefore, that matter is substance. But this is impossible. For separateness and that ``something'' seem to belong to substance. \\
\hline
\textbf{Final Definition} & Let us now look at the residue (\textit{Residuum}) of the products of labour. There is nothing left of them in each case but the same ghost-like (\textit{gespenstige}) objectivity; they are merely gelatin (\textit{Gallerte}) of homogeneous human labour, i.e.\ of human labour-power expended without regard to the form of its expenditure. \dots\ As crystals (\textit{Kristalle}) of this social substance, which is common to them all, they are values---commodity values. & [above] \\
%\hline
\end{tabularx}
\end{table}
\end{singlespacing}

\newpage
\section{Lysis}
Theoretical consequences hold, though, regardless of whether Marx explicitly referenced the passage. It's relatively clear, then, why Aristotle can only fail, in \textit{The Nicomachean Ethics}, to come up with congealed labor-time from his analysis of exchange. Not only would such an endpoint be outside the scope of that ethical discussion, but such a concept would strictly contradict Aristotle's understanding of this sort of universalized predicate as something that can maintain neither quantity nor quality---even were Aristotle to limit the domain to his own ethical-exchange discussion.\footnote{In fact Aristotle almost certainly would not have combined these two domains in the first place. Aristotle wouldn't have collapsed his metaphysical system with his ethical system, since he would have conceived of them as relatively separate modes of inquiry (perhaps displaying a truer effect of Aristotle's era constraints). Therein, my own smashing together of the \textit{Ethics} and \textit{Metaphysics} as a justification for a certain argument is, in a certain light, not germane. My smashing could only follow Marx's lead in collapsing the ontological and ethical, even if this particular smashing of his appears more false than my own. I owe this insight to Cesare Casarino.} Directly influenced or not, Marx has redefined, assuming Aristotle's perspective, the homogenous as yet quantitatively and qualitatively particular. By this redefinition, Marx's blustery translation of the \textit{Ethics} can be, as I promised, restituted.  

The \textit{Grundrisse}, in fact, slides open a window for understanding Marx's mistranslation in \textit{Capital}. Marx provides there a sardonic snippet from the exact same passage of the \textit{Ethics} as the one I've been discussing here from \textit{Capital}. ``Money is a dead pledge''\footnote{Marx, \textit{Grundrisse}, 160.} he quotes, which is a construal of what more literally would be translated as, ``As for the expectation of exchange, money is our guarantor; even if there is no need (for it) now, because it will exist if it becomes needed.''\footnote{Aristotle, \textit{Ethics}, 286.} Following the quote, Marx explicitly claims that this exchange relation is a form of an individual's ``alienation (of) their own social relationship from themselves''.\footnote{Marx, \textit{Grundrisse}, 160.} It seems, then, that Marx might be thinking of this particular passage, among the many others, in terms of the alienation of humans under a network of commodity exchange.  

Why Marx gratuitously and disingenuously paraphrased Aristotle's ``commensurability'' as a ``qualitative equality'' and a ``substance'' can be thrown into further relief, now. Marx's gelatinous labor can only maintain its status as ``alienated'' if it independently maintains some quantity and quality of its own, if it is in some sense particular. Were this not the case, the commodity domain in which Marx delimits his universal predicate, would collapse into Aristotle's own unlimited domain. The alienation of labor, which I must now define as alienation between domains, would fail to gain conceptual purchase.\footnote{The argument is that if the domain of the particularized commodity is non-unique, workers and commodities would be de facto part of the same set of entities, in which case it would be difficult to make a claim that commodities are alienated. Of course, Marx does characterize labor-power as a commodity, and also workers as commodities in a certain sense. But workers and their labor-power, I think, would still need to maintain at least formal membership in distinct sets---similar to the not-formal but \textit{de facto} slavery of free wage-laborers.}  

Marx, then, inserts value qua substance, into and against Aristotle's text, as a way of highlighting that labor substance's very detachableness, i.e. its alienating power, its ability to maintain itself. It's not the case that Aristotle's historical situation prevented him from gaining insight into the nature of congealed labor time, but rather that the alienated nature of the commodity-value can assert itself over and against another historical account of the process of exchange---though this self-assertion, I'll add, is still dependent on the commodity being society's universal form.\footnote{Thus there is, again, still societal constraints. My reading here can maintain a spectrum of implications, based on Marx's level of  awareness, of which I'll characterize the endpoints. On one extreme, Marx is completely unaware of his insertion of the commodity substance, and I read Marx's own claim against himself---his own circumstances in a commodity-based society blind him to his ahistorical commodity-assertion, and he fails to fully account for historical progression. On the other extreme, Marx is completely aware of the insertion, and I read Marx as highlighting the extreme corrosiveness of the commodity-world.}

\section{Conclusion}
Finally, I'll commit to and end with my own sin of bluster. I'll note that congealed labor is alienating in another sense. \textit{Gallerte}, that mass of crushed saints, would likely have been eaten by the wide-eyed readers of \textit{Capital}, forcing them to reckon with the image of a quite alienated laborer's bones squeaking against their teeth---Marx meant it to horrify.\footnote{Ngai, \textit{Categories}, 67.} Thus, bigboy Marx, reprising his dig against all those who failed in the previous 2000 years to understand the ``slight in content''\footnote{Marx, \textit{Capital}, 90.} law-of-equivalence, commits the first recorded, to my knowledge, act of forced textual gelatophagy against the Philosopher-capital-P, by metaphorically forcing that labor-owning, aristocratic Aristotle to swallow the residue from capital's grinding obsolete of Aristotle's concept of exchange.

\section{References}
%\addcontentsline{toc}{section}{Works Cited}
\begin{singlespacing}
\begin{hangparas}{.25in}{1}
Aristotle. \textit{The Metaphysics I--IX}. Translated by Hugh Tredennick. Loeb Classical Library 271. Cambridge, MA: Harvard University Press, 1933.

---. \textit{The Nicomachean Ethics}. Translated by H. Rackham. Loeb Classical Library 73. Cambridge, MA: Harvard University Press, 1934.

---. \textit{Politics}. Translated by H. Rackham. Loeb Classical Library 264. Cambridge, MA: Harvard University Press, 1934.

Derrida, Jacques. \textit{Specters of Marx: The State of the Debt, the Work of Mourning and the New International}. Translated by Peggy Kamuf. New York: Routledge, 1994.

Heidegger, Martin. ``Letter on Humanism''. In \textit{Basic Writings}. Edited by David Farrell Krell. San Francisco: Harper, 1993.

Marx, Karl. \textit{Capital Volume I}. Translated by Ben Fowkes. New York: Penguin Books, 1990.

---. \textit{The Difference Between the Democritean and Epicurean Philosophy of Nature}. 1841.

---. \textit{Grundrisse}. Translated by Martin Nicolaus. New York: Penguin Books, 1990.

Ngai, Sianne. \textit{Our Aesthetic Categories: Zany, Cute, Interesting}. Cambridge, MA: Harvard University Press, 2012.

Postone, Moishe. \textit{Time, Labor, and Social Domination: A Reinterpretation of Marx's Critical Theory}. Cambridge: Cambridge University Press, 1993.
\end{hangparas}
\end{singlespacing}

\end{document}
